\documentclass[oneside,numbers,spanish,a4paper]{ezthesis/ezthesis}
%% # Opciones disponibles para el documento #
%%
%% Las opciones con un (*) son las opciones predeterminadas.
%%
%% Modo de compilar:
%%   draft            - borrador con marcas de fecha y sin im'agenes
%%   draftmarks       - borrador con marcas de fecha y con im'agenes
%%   final (*)        - version final de la tesis
%%
%% Tama'no de papel:
%%   letterpaper (*)  - tama'no carta (Am'erica)
%%   a4paper          - tama'no A4    (Europa)
%%
%% Formato de impresi'on:
%%   oneside          - hojas impresas por un solo lado
%%   twoside (*)      - hojas impresas por ambos lados
%%
%% Tama'no de letra:
%%   10pt, 11pt, o 12pt (*)
%%
%% Espaciado entre renglones:
%%   singlespace      - espacio sencillo
%%   onehalfspace (*) - espacio de 1.5
%%   doublespace      - a doble espacio
%%
%% Formato de las referencias bibliogr'aficas:
%%   numbers          - numeradas, p.e. [1]
%%   authoryear (*)   - por autor y a'no, p.e. (Newton, 1997)
%%
%% Opciones adicionales:
%%   spanish         - tesis escrita en espa'nol
%%
%% Desactivar opciones especiales:
%%   nobibtoc   - no incluir la bibiolgraf'ia en el 'Indice general
%%   nofancyhdr - no incluir "fancyhdr" para producir los encabezados
%%   nocolors   - no incluir "xcolor" para producir ligas con colores
%%   nographicx - no incluir "graphicx" para insertar gr'aficos
%%   nonatbib   - no incluir "natbib" para administrar la bibliograf'ia

%% Paquetes adicionales requeridos se pueden agregar tambi'en aqu'i.
%% Por ejemplo:
%\usepackage{subfig}
%\usepackage{multirow}
\usepackage[utf8]{inputenc}
\usepackage{mathptmx} % Times new roman
\usepackage{enumitem}
\usepackage[export]{adjustbox}
\usepackage[figurename=Fig.,tablename=TABLA]{caption}
\usepackage{float}
\usepackage{flafter}
\usepackage{afterpage}
\usepackage{placeins}
\usepackage{tabularx}
\usepackage{chngcntr}
\usepackage{titlesec}
\usepackage{lscape}
\usepackage{listings} % code
\usepackage{color}
\usepackage{textcomp} % quotesimple
\usepackage{multicol}

\definecolor{codegreen}{rgb}{0,0.6,0}
\definecolor{codegray}{rgb}{0.5,0.5,0.5}
\definecolor{codepurple}{rgb}{0.58,0,0.82}
\definecolor{backcolour}{rgb}{0.95,0.95,0.92}
 
\lstdefinestyle{codestyle}{
    backgroundcolor=\color{backcolour},
    commentstyle=\color{codegreen},
    keywordstyle=\color{magenta},
    numberstyle=\tiny\color{codegray},
    stringstyle=\color{codepurple},
    basicstyle=\footnotesize,
    breakatwhitespace=false,
    breaklines=true,
    captionpos=b,
    keepspaces=false,
    numbers=left,
    numbersep=5pt,
    showspaces=false,
    showstringspaces=false,
    showtabs=false,
    tabsize=2,
    inputencoding=utf8,
    literate={á}{{\'a}}1 {é}{{\'e}}1 {í}{{\'i}}1 {ó}{{\'o}}1 {ú}{{\'u}}1
}
\lstset{style=codestyle}

%% # Datos del documento #
%% Nota que los acentos se deben escribir: \'a, \'e, \'i, etc.
%% La letra n con tilde es: \~n.

\author{Kelvin Helmut Provincia Quispe}
\title{PLANTILLA DE PROYECTO DE FIN DE CARRERA}
\degree{[Grado o Título a obtener]}
\supervisor{Nombre de mi Asesor}
\institution{Universidad La Salle}
\faculty{ESCUELA PROFESIONAL DE INGENIERÍA DE SOFTWARE}
\department{AREQUIPA – PERÚ}

%% # M'argenes del documento #
%% 
%% Quitar el comentario en la siguiente linea para ajustar los m'argenes del
%% documento. Leer la documentaci'on de "geometry" para m'as informaci'on.

%\geometry{top=40mm,bottom=33mm,inner=40mm,outer=25mm}

%% El siguiente comando agrega ligas activas en el documento para las
%% referencias cruzadas y citas bibliogr'aficas. Tiene que ser *la 'ultima*
%% instrucci'on antes de \begin{document}.

\hyperlinking
\begin{document}

%% En esta secci'on se describe la estructura del documento de la tesis.
%% Consulta los reglamentos de tu universidad para determinar el orden
%% y la cantidad de secciones que debes de incluir.


\counterwithout{figure}{chapter} % chngcntr
\counterwithout{table}{chapter}

\renewcommand\thechapter{\Roman{chapter}}
\renewcommand{\thesubsection}{\Alph{subsection}}
\renewcommand{\thesection}{\arabic{chapter}.\arabic{section}}
\renewcommand\listtablename{Índice de tablas}
\renewcommand\bibname{Referencias}
\renewcommand\thefigure{\arabic{figure}}
\renewcommand\thetable{\Roman{table}}

\titleformat{\chapter}[hang] {\vspace{-6em}\centering\normalfont\LARGE\bfseries}{\chaptertitlename\ \thechapter\ --}{0.2em}{}

\captionsetup[table]{
  labelsep=newline,
  justification=centering,
}

\fancyfoot[R]{\thepage}
\fancyhead[R]{}

%% # Portada de la tesis #
%% Mirar el archivo "titlepage.tex" para los detalles.
%% ## Construye tu propia portada ##
%% 
%% Una portada se conforma por una secuencia de "Blocks" que incluyen
%% piezas individuales de informaci'on. Un "Block" puede incluir, por
%% ejemplo, el t'itulo del documento, una im'agen (logotipo de la universidad),
%% el nombre del autor, nombre del supervisor, u cualquier otra pieza de
%% informaci'on.
%%
%% Cada "Block" aparece centrado horizontalmente en la p'agina y,
%% verticalmente, todos los "Blocks" se distruyen de manera uniforme 
%% a lo largo de p'agina.
%%
%% Nota tambi'en que, dentro de un mismo "Block" se pueden cortar
%% lineas usando el comando \\
%%
%% El tama'no del texto dentro de un "Block" se puede modificar usando uno de
%% los comandos:
%%   \small      \LARGE
%%   \large      \huge
%%   \Large      \Huge
%%
%% Y el tipo de letra se puede modificar usando:
%%   \bfseries - negritas
%%   \itshape  - it'alicas
%%   \scshape  - small caps
%%   \slshape  - slanted
%%   \sffamily - sans serif
%%
%% Para producir plantillas generales, la informaci'on que ha sido inclu'ida
%% en el archivo principal "tesis.tex" se puede accesar aqu'i usando:
%%   \insertauthor
%%   \inserttitle
%%   \insertsupervisor
%%   \insertinstitution
%%   \insertdegree
%%   \insertfaculty
%%   \insertdepartment
%%   \insertsubmitdate

\begin{titlepage}
  \TitleBlock[\vspace{-4cm}\hspace{-0.8cm}]{\includegraphics[width=6cm,left]{uls}}
%  \TitleBlock{\scshape\insertinstitution}
  \TitleBlock[\vspace{2cm}]{\scshape\insertfaculty}
  \TitleBlock[\vspace{1.5cm}]{\large\scshape\inserttitle}
  \TitleBlock[\vspace{1.5cm}]{\scshape\insertauthor}
  \TitleBlock[\bigskip\bigskip]{\scshape\insertdegree}
  \TitleBlock{\insertdepartment}
  \TitleBlock[\bigskip]{\insertsubmitdate}
\end{titlepage}

%% Nota 1:
%% Se puede agregar un escudo o logotipo en un "Block" como:
%%   \TitleBlock{\includegraphics[height=4cm]{escudo_uni}}
%% y teniendo un archivo "escudo_uni.pdf", "escudo_uni.png" o "escudo_uni.jpg"
%% en alg'un lugar donde LaTeX lo pueda encontrar.

%% Nota 2:
%% Normalmente, el espacio entre "Blocks" se extiende de modo que el
%% contenido se reparte uniformemente sobre toda la p'agina. Este
%% comportamiento se puede modificar para mantener fijo, por ejemplo, el
%% espacio entre un par de "Blocks". Escribiendo:
%%   \TitleBlock{Bloque 1}
%%   \TitleBlock[\bigskip]{Bloque2}
%% se deja un espacio "grande" y de tama~no fijo entre el bloque 1 y 2.
%% Adem'as de \bigskip est'an tambi'en \smallskip y \medskip. Si necesitas
%% aun m'as control puedes usar tambi'en, por ejemplo, \vspace*{2cm}.




%% # Prefacios #
%% Por cada prefacio (p.e. agradecimientos, resumen, etc.) crear
%% un nuevo archivo e incluirlo aqu'i.
%% Para m'as detalles y un ejemplo mirar el archivo "gracias.tex".
% \prefacesection{Dedicatoria}
\thispagestyle{empty}\mbox{}
\vspace*{4cm}
\begin{flushright}
    \emph{Dedicatoria}
\end{flushright}
%% Las secciones del "prefacio" inician con el comando \prefacesection{T'itulo}
%% Este tipo de secciones *no* van numeradas, pero s'i aparecen en el 'indice.
%%
%% Si quieres agregar una secci'on que no vaya n'umerada y que *tampoco*
%% aparesca en el 'indice, usa entonces el comando \chapter*{T'itulo}
%%
%% Recuerda que aqu'i ya puedes escribir acentos como: 'a, 'e, 'i, etc.
%% La letra n con tilde es: 'n.

\prefacesection{Agradecimientos}


%% Por si alguien tiene curiosidad, este "simp'atico" agradecimiento est'a
%% tomado de la "Tesis de Lydia Chalmers" basada en el universo del programa
%% de televisi'on Buffy, la Cazadora de Vampiros.
%% http://www.buffy-cazavampiros.com/Spiketesis/tesis.inicio.htm


%% # 'Indices y listas de contenido #
%% Quitar los comentarios en las lineas siguientes para obtener listas de
%% figuras y cuadros/tablas.
\tableofcontents
\prefacesection{Índice de abreviaturas y siglas}

\listoftables
\listoffigures
\prefacesection{Glosario de términos}

\prefacesection{Resumen}
Resumen\cite{NewCam97}

\prefacesection{Abstract}
\prefacesection{Palabras clave}


%% # Cap'itulos #
%% Por cada cap'itulo hay que crear un nuevo archivo e incluirlo aqu'i.
%% Mirar el archivo "intro.tex" para un ejemplo y recomendaciones para
%% escribir.
\chapter{Problemática del Proyecto}

\section{Contexto del problema}

\section{Antecedentes}

\chapter{Planteamiento del Proyecto}
\section{Fundamentos teóricos}

\section{Objetivos del Proyecto}
\subsection{Objetivo General}

\subsection{Objetivos Específicos}

\section{Justificación}

\section{Viabilidad}
\subsection{Viabilidad técnica}

\subsection{Viabilidad operativa}

\subsection{Viabilidad económica}

\section{Limitaciones}

\chapter{Metodología}
%\include{conclu}

\appendix
%% Cap'itulos incluidos despues del comando \appendix aparecen como ap'endices
%% de la tesis.
%\include{apendiceA}
%\include{apendiceB}
%\include{apendiceC}

%% Incluir la bibliograf'ia. Mirar el archivo "biblio.bib" para m'as detales
%% y un ejemplo.
% Ver ezthesis.cls linea 178
%\bibliographystyle{IEEEannot}
\bibliography{bibliografia/biblio}

\end{document}
